%%%%%%%%%%%%%%%%%%%%%%%%%%%%%%%%%%%%%%%%%%%%%%%%%%%%%%%%%%%%%%%%%%%%%%%%

%%% LaTeX Template for AAMAS-2026 (based on sample-sigconf.tex)
%%% Prepared by the AAMAS-2026 Publication Chairs based on the version from AAMAS-2025. 

%%%%%%%%%%%%%%%%%%%%%%%%%%%%%%%%%%%%%%%%%%%%%%%%%%%%%%%%%%%%%%%%%%%%%%%%

%%% Start your document with the \documentclass command.

\documentclass[sigconf,anonymous]{aamas} 

%%% Load required packages here (note that many are included already).

\usepackage{balance} % for balancing columns on the final page
\usepackage{algorithm}
\usepackage{algorithmic}
\usepackage{booktabs}
\usepackage{multirow}
\usepackage{colortbl}
\usepackage{xcolor}
\usepackage[section]{placeins}
\usepackage{amsmath}
\DeclareMathOperator*{\argmax}{arg\,max}
\DeclareMathOperator*{\argmin}{arg\,min}
\usepackage{caption}
\usepackage{subcaption}
\usepackage{todonotes}
\usepackage{orcidlink}
\usepackage{hyperref}


%%%%%%%%%%%%%%%%%%%%%%%%%%%%%%%%%%%%%%%%%%%%%%%%%%%%%%%%%%%%%%%%%%%%%%%%

%%% AAMAS-2026 copyright block (do not change!)

\setcopyright{ifaamas}
\acmConference[AAMAS '26]{Proc.\@ of the 25th International Conference
on Autonomous Agents and Multiagent Systems (AAMAS 2026)}{May 25 -- 29, 2026}
{Paphos, Cyprus}{C.~Amato, L.~Dennis, V.~Mascardi, J.~Thangarajah (eds.)}
\copyrightyear{2026}
\acmYear{2026}
\acmDOI{}
\acmPrice{}
\acmISBN{}


%%%%%%%%%%%%%%%%%%%%%%%%%%%%%%%%%%%%%%%%%%%%%%%%%%%%%%%%%%%%%%%%%%%%%%%%

%%% == IMPORTANT ==
%%% Use this command to specify your submission number.
%%% In anonymous mode, it will be printed on the first page.

\acmSubmissionID{<<submission id>>}

%%% Use this command to specify the title of your paper.

\title[Counterfactual Gradient Alignment]{Counterfactual Gradient Alignment: Enhancing Data Efficiency through Directional Expert Supervision}

%%% Provide names, affiliations, and email addresses for all authors.

\author{Arthur Pendragon}
\affiliation{
  \institution{Camelot Castle}
  \city{Camelot}
  \country{United Kingdom}}
\email{king.arthur@camelot.uk}

\author{Nimue}
\affiliation{
  \institution{The Lady's Lake}
  \city{Avalon}
  \country{United Kingdom}}
\email{lady.of.the.lake@avalon.uk}

%%% Use this environment to specify a short abstract for your paper.

\begin{abstract}
Supervised learning typically relies on passive data-label pairs, often failing to capture the rich geometric intuition that human experts possess about decision boundaries. This limitation is particularly acute in low-data regimes where models lack sufficient signal to generalize. We propose \textbf{Counterfactual Gradient Alignment (CGA)}, a framework that leverages expert-provided directions pointing towards counterfactual instances to supervise the model's gradient landscape. By penalizing positive directional derivatives along these vectors, we effectively teach the model where class transitions should occur. We evaluate three novel loss formulations---ReLU, Softplus, and Sign-based variants---across a diverse suite of benchmarks including synthetic 2D datasets, image classification (MNIST), and sentiment analysis (IMDB, SST). Our results demonstrate significant accuracy boosts, such as a +7.1\% gain on IMDB with only 200 samples and a +14.5\% improvement on complex 2D boundaries. Crucially, we identify a profound synergy between CGA and active learning, demonstrating that directional supervision is most impactful when applied to model-selected informative boundary samples.
\end{abstract}

%%% Use this command to specify a few keywords describing your work.
%%% Keywords should be separated by commas.

\keywords{Counterfactuals, Gradient Supervision, Active Learning, Human-AI Interaction}

%%%%%%%%%%%%%%%%%%%%%%%%%%%%%%%%%%%%%%%%%%%%%%%%%%%%%%%%%%%%%%%%%%%%%%%%

%%% Include any author-defined commands here.         
\newcommand{\BibTeX}{\rm B\kern-.05em{\sc i\kern-.025em b}\kern-.08em\TeX}

%%%%%%%%%%%%%%%%%%%%%%%%%%%%%%%%%%%%%%%%%%%%%%%%%%%%%%%%%%%%%%%%%%%%%%%%

\begin{document}

%%% The following commands remove the headers in your paper. For final 
%%% papers, these will be inserted during the pagination process.

\pagestyle{fancy}
\fancyhead{}

%%% The next command prints the information defined in the preamble.

\maketitle 

%%%%%%%%%%%%%%%%%%%%%%%%%%%%%%%%%%%%%%%%%%%%%%%%%%%%%%%%%%%%%%%%%%%%%%%%

\section{Introduction}
Modern machine learning models, while powerful, often require vast amounts of data to learn reliable decision boundaries. In contrast, humans can generalize from a handful of examples by applying structural priors and understanding class relationships. A human expert doesn't just know that an image is a `3'; they know that changing a specific curve would turn it into an `8'. This directional intuition represents a missed opportunity for model supervision.

We introduce \textbf{Counterfactual Gradient Alignment (CGA)}, which converts this intuitive knowledge into a differentiable training signal. In CGA, an expert provides a direction vector $\mathbf{d}$ indicating the shortest path to a counterfactual (a point where the classification changes). We then constrain the model's gradients to align with this vector, ensuring that moving towards the counterfactual consistently reduces the model's confidence in the original class.

This paper makes three primary contributions: (1) we define a family of directional loss functions (ReLU, Softplus, Sign) for gradient supervision; (2) we demonstrate that CGA significantly improves generalization in low-data regimes across vision and NLP; and (3) we show that CGA is uniquely dependent on the informativeness of the training samples, thriving in active learning settings while struggling with random data subsets.

\section{Related Work}
Our work builds on \textbf{Gradient Supervision} \cite{teney2020learning}, which penalizes the angle between model gradients and counterfactual vectors. While Teney et al. focus on angular alignment, our approach explores magnitude-aware penalties like Softplus and bounded signals like Sign/Tanh. We also draw from \textbf{Counterfactually Augmented Data} \cite{Kaushik2020Learning}, which uses human-edited reviews to improve NLP robustness. Unlike Kaushik et al., who use counterfactuals as additional data points, we use them to supervise the model's local derivative.

\section{Proposed Method: CGA}
Given an input $\mathbf{x}$ and its counterfactual $\hat{\mathbf{x}}$, the \textbf{Directional Derivative} $d$ is:
begin{equation}
    d = \nabla_x f_{y}(x) \cdot \frac{\hat{\mathbf{x}} - \mathbf{x}}{||\hat{\mathbf{x}} - \mathbf{x}||}
end{equation}
where $f_y(x)$ is the model's score for the true class $y$. A negative $d$ indicates desired behavior (confidence drops towards the boundary). We minimize $d$ using:
\begin{itemize}
    \item \textbf{Softplus CGA}: $\mathcal{L}_{\text{SP}} = \log(1 + e^{\beta d})$. This smooth variant encourages a continuous drop in confidence.
    \item \textbf{Sign CGA}: $\mathcal{L}_{\text{Sign}} = \tanh(\beta d) + 1$. This variant squashes outliers and focuses on the sign of the derivative.
\end{itemize}
The total loss is: $\mathcal{L}_{\text{Total}} = (1-\alpha)\mathcal{L}_{\text{CE}} + \alpha\mathcal{L}_{\text{CGA}}$.

\section{Experiments and Results}

\subsection{Experimental Setup}
For vision, we use a CNN architecture on MNIST. For NLP, we use a Bag-of-Words neural network with trainable embeddings on IMDB, SST-2, and SST-5. We simulate low-data scenarios by subsetting the training data to as few as 10 samples.

\subsection{Text Classification Results}
CGA achieves substantial gains on IMDB and SST-5 when paired with active learning (AL). As shown in Table \ref{tab:nlp_results}, Softplus is the most consistent performer.

\begin{table}[h]
\centering
\caption{Validation Accuracy on IMDB and SST-5 (AL ON).}
\label{tab:nlp_results}
\begin{tabular}{lcccc}
\toprule
\textbf{Dataset} & \textbf{Size} & \textbf{Baseline} & \textbf{CGA (Ours)} & \textbf{Gain} \\ \midrule
IMDB & 200 & 55.9\% & \textbf{63.1\% (SP)} & \textbf{+7.2\%} \\
IMDB & 100 & 51.4\% & \textbf{56.6\% (Sign)} & \textbf{+5.2\%} \\
SST-5 & 100 & 53.1\% & \textbf{56.5\% (ReLU)} & \textbf{+3.4\%} \\
\bottomrule
\end{tabular}
\end{table}

\subsection{Active Learning Synergy}
We observed that CGA's performance flips from positive to negative when AL is disabled (Table \ref{tab:al_synergy}). This suggests that directional supervision on "easy" (central) samples provides little value or introduces noise, whereas supervising boundary samples provides a strong regularizing signal.

\begin{table}[h]
\centering
\caption{Performance Gain over Baseline (Size 100) with/without Active Learning.}
\label{tab:al_synergy}
\begin{tabular}{lcc}
\toprule
\textbf{Condition} & \textbf{AL Enabled} & \textbf{AL Disabled} \\ \midrule
Gain (IMDB) & \textbf{+5.1\%} & -2.3\% \\
Gain (SST-5) & \textbf{+3.4\%} & -4.8\% \\
\bottomrule
\end{tabular}
\end{table}

\section{Discussion and Conclusion}
CGA demonstrates that model gradients can be explicitly supervised using counterfactual directions to achieve superior generalization in data-scarce environments. The high performance of the \textbf{Softplus} variant suggests that smooth, persistent gradient alignment is preferable to hard hinge penalties. Future work will explore the scalability of CGA to large transformer architectures and its integration into interactive human-AI teaching loops.

\balance
\bibliographystyle{ACM-Reference-Format}
\bibliography{sample}

\end{document}
%%%%%%%%%%%%%%%%%%%%%%%%%%%%%%%%%%%%%%%%%%%%%%%%%%%%%%%%%%%%%%%%%%%%%%%%
